% New proof: central sign translates and positive squared dimensions.
% This file is independent of any unproved character/monodromy identification.
\section{Positive central moments and the sign object}
\label{sec:frontier-signed-gauss}

We use the canonical spherical structure on the sphericalization
$\mathcal S$ of a hypothetical fusion category. Its realization as an
equivariantization for the canonical involutive pivotal functor gives
a central copy of $\operatorname{Rep}(C_2)$; see \cite[Section~5.3]{Bartlett}. Its central sign object
$\varepsilon$ has order two, dimension $-1$, and self-braiding $+1$.
Consequently its twist in $\mathcal Z(\mathcal S)$ is $-1$.
The two lifts of each original self-dual simple are self-dual; their
dimensions and second Frobenius--Schur indicators have opposite signs.
The latter assertion follows directly from the definition of the second
indicator: the two lifts have the same evaluation and coevaluation maps
after forgetting equivariance, while their pivotal maps differ by $-1$.

\begin{lemma}[A positive measure without pseudo-unitarity]
\label{lem:positive-central-squares}
Let $\mathcal S$ be a spherical fusion category, and put
$d=\dim\mathcal S$. For $A\in\operatorname{Irr}\mathcal Z(\mathcal S)$,
let $a_A=\dim A$ and $\theta_A$ be its twist. Then $a_A$ is real and
\[
 \sum_A a_A^2=d^2,\qquad
 \sum_A a_A^2\theta_A=d,\qquad
 \sum_A a_A^2\theta_A^2=d\sum_V\dim(V)\nu_2(V).
\]
In particular the coefficients $a_A^2$ form a positive measure on roots
of unity, even when some dimensions of $\mathcal S$ are negative.
\end{lemma}
\begin{proof}
Simple dimensions in a spherical fusion category are totally real
\cite[Section~2.1]{OstrikRank3}. The forgetful functor is spherical, so
$a_A$ is an integral linear combination of these real dimensions. Let $I$ be
central induction and $m_{V,A}=[I(V):A]$. Adjunction gives
$a_A=\sum_Vm_{V,A}\dim V$.
The indicator identity of \cite[Theorem~4.1]{NgSchauenburg}
\[
 \operatorname{Tr}(\theta^n\vert_{I(V)})=d\nu_n(V)
\]
therefore gives, after multiplication by $\dim V$ and summation over $V$,
$\sum_A a_A^2\theta_A^n=d\sum_V\dim(V)\nu_n(V)$.
For $n=1$ only the tensor unit contributes, and for $n=2$ this is the
claimed formula. The zeroth moment is the standard identity
$\dim\mathcal Z(\mathcal S)=d^2$. The twist is a root of unity by the
Vafa theorem for the center of a spherical fusion category.
\end{proof}

\begin{proposition}[A global codegree moment bound]\label{prop:spherical-codegree-moment}
Let a commutative fusion ring have a spherical categorification, with
simple dimensions $d_i$, global dimension $\Delta=\sum_i d_i^2$ and
formal codegrees $f_j$, counted with multiplicity. Put
\[
 w=\sum_j f_j^{-2},\qquad
 S=\sum_{i=i^*}|d_i|,\qquad h=\Delta^2w.
\]
At every real embedding of its dimension field one has
\begin{equation}\label{eq:spherical-codegree-moment}
 2(\Delta-h)^2\leq
 (\Delta^2-h)(\Delta^2+\Delta S-2h).
\end{equation}
Equivalently,
$\Delta(1-3w)+(1-w)S+4w-2/\Delta\geq0$.
For a self-dual pseudo-unitary category,
$S=\sum_i\operatorname{FPdim}(X_i)$.
\end{proposition}
\begin{proof}
The constituents of $I(1)$ are distinct, have twist one and dimensions
$\Delta/f_j$. Remove their squared mass $h$ from the positive measure
of Lemma~\ref{lem:positive-central-squares}. The remaining mass is
$\Delta^2-h$, its first moment is $\Delta-h$, and its second moment is
$\Delta s-h$, where $s=\sum_i d_i\nu_2(X_i)\leq S$.
The weighted sum of squared real parts of twists is
$(\Delta^2+\Delta s-2h)/2$. Cauchy--Schwarz gives the displayed
inequality with $s$ in place of $S$; increasing $s$ to $S$ preserves it
because the remaining mass is nonnegative. Expanding and dividing by
$\Delta^3>0$ gives the equivalent expression.
The argument applies to every Galois conjugate spherical structure.
The rational number $w$ can be computed without finding codegree roots:
$w=\operatorname{Tr}(B^{-2})$ for the Casimir matrix $B$.
\end{proof}

This form of the bound applies to every possible spherical dimension
character, including those outside the Frobenius--Perron orbit. In the
finite certificates it excludes the remaining character branches for
$F_{21,325}$ and $F_{24,735}$; a pseudo-unitary obstruction alone would
not exclude those branches.

\begin{proposition}[The signed-translate moment obstruction]
\label{prop:signed-translate-moment}
Let $\mathcal S$ be a sphericalization with
$\dim\mathcal S=2\delta$ and commutative Grothendieck ring. Write
$A_j$ for the simple constituents of $I(1)$ and $a_j=\dim A_j>0$.
Put
\[
 h=\sum_j a_j^2,\qquad
 2s=\sum_{V\in\operatorname{Irr}\mathcal S}\dim(V)\nu_2(V).
\]
Then there are signs $\eta_j\in\{1,-1\}$, with $\eta_j=-1$ for
$A_j=1$, such that
\[
 \sum_j\eta_ja_j=0.
\]
For $C=\sum_j\eta_ja_j^2$, one necessarily has
\begin{gather*}
 h\leq2\delta^2,\qquad h/\delta-\delta\leq s\leq\delta,\\
 (2\delta-h-C)^2
 \leq(4\delta^2-2h)(2\delta s+2\delta^2-2h).
\end{gather*}
These conditions hold after every Galois conjugation, with the same
second indicators and twist signs.
\end{proposition}
\begin{proof}
The formal-codegree theorem \cite[Theorem~2.13]{OstrikRank3} shows that
$I(1)$ is multiplicity free and its constituents have positive dimensions;
their twists are one by \cite[Theorem~2.5]{OstrikRank3}. The projection formula gives
$I(\varepsilon)=\varepsilon\otimes I(1)$, whose constituents have
dimensions $-a_j$ and are therefore disjoint from those of $I(1)$.
Their twists $\eta_j$ are $\pm1$, by balancing and the fact that the
monodromy with an invertible object of order two squares to one.
The indicator identity for $I(\varepsilon)$ gives
$\sum_j\eta_ja_j=0$; the constituent $\varepsilon\otimes1$ has twist
$-1$.

Remove these two sets of constituents from the positive measure of
Lemma~\ref{lem:positive-central-squares}. The remaining mass is
$M=4\delta^2-2h$, its first moment is $2\delta-h-C$, and its second
moment is $4\delta s-2h$. If its twists are $e^{i\phi_A}$, then
\[
 \sum_{\rm remaining}a_A^2\cos^2\phi_A
 =2\delta s+2\delta^2-2h.
\]
The inequalities $0\leq\sum a_A^2\cos^2\phi_A\leq M$ give the
two bounds on $s$. Cauchy--Schwarz gives the last inequality. Every
dimension remains real under the Galois embeddings under consideration,
whereas the second indicators and the twist signs are rational numbers
and remain fixed.
\end{proof}


\begin{corollary}[Finite signed-moment certificate]\label{cor:sign-certificate}
Let $T$ be a proposed signed table from Proposition~\ref{prop:signed}, and
let $\chi$ be a proposed dimension character, with
$\delta=\sum_i\chi(b_i)^2$. Form the eight numbers $a_j=\delta/f_j$,
where the $f_j$ are the formal codegrees of $R$ and $T$, with all
multiplicities. Choose the occurrence $f_j=\delta$ corresponding to the
central unit. For every self-dual nonunit basis element choose
$\nu_i\in\{-1,1\}$ and put
\[
 s=1+\sum_{i>0,\ i=i^*}\nu_i\chi(b_i).
\]
If no joint choice of these indicators and the twist signs $\eta_j$
satisfies Proposition~\ref{prop:signed-translate-moment} at every real
embedding, $T$ cannot be the signed quotient of a sphericalization with
dimension character $\chi$.
\end{corollary}
\begin{proof}
The two lifts of a non-self-dual object have zero second indicator.
For a self-dual object the two lift dimensions and indicators both change
sign, so their products agree. The unit and its sign translate contribute
two. This gives the displayed $s$ and exhausts the indicator possibilities.
The formal-codegree description in Proposition~\ref{prop:signed} gives the
eight $a_j$. The remaining assertions are exactly the necessary conditions
of Proposition~\ref{prop:signed-translate-moment}. Equal codegrees retain
separate occurrences; an occurrence of dimension one may be designated the
unit without losing any possibility, since the other signs are all enumerated.
\end{proof}

The supplied certificate for $F_{17,308}$ uses this corollary in one
character branch. The other signed branches are handled by the same
arithmetic, spherical and parity routines as all other rings.
